\chapterstar{Justificación}

El jitomate es uno de esos cultivos que sostienen a miles de familias en México. No es solo un negocio: para muchos productores representa el ingreso con el que pagan deudas, mandan a sus hijos a la escuela y planean el futuro. Por eso, cuando una plaga o una enfermedad llega sin avisar, el golpe no es solo económico, es personal.

Hoy en día, la mayoría de los productores todavía tienen que salir ellos mismos a caminar entre las plantas, fila por fila, bajo el sol, buscando señales de daño. Es un trabajo agotador y, en terrenos grandes, prácticamente imposible de hacer bien. Para cuando se nota que algo anda mal, muchas veces ya se propagó demasiado.

Ahí es donde nace la idea de este proyecto. En lugar de esperar a que el productor camine hasta el problema, la propuesta es desarrollar un sistema de monitoreo que usa un dron equipado con cámara y GPS para sobrevolar las áreas de cultivo en rutas programadas. Mientras el dron hace su recorrido, va capturando imágenes de las plantas de jitomate; regresa a su punto de inicio y las envía a través de WiFi, que proporciona un router con un chip 4G, a la Raspberry Pi~4.

Ahí entra la inteligencia artificial: se utilizarán herramientas de visión artificial e inteligencia artificial capaces de analizar imágenes capturadas por el dron de manera automatizada. Entre las principales tecnologías implementadas se encuentran \textbf{OpenCV}, cuya función principal será recibir y procesar las fotografías capturadas por el dron antes de ser analizadas por la inteligencia artificial, y \textbf{TensorFlow Lite}, cuya función dentro del proyecto será ejecutar el modelo de red neuronal encargado de analizar las imágenes y detectar posibles anomalías en los cultivos de jitomate.

El modelo de inteligencia artificial será entrenado previamente con imágenes de plantas sanas y plantas afectadas por enfermedades o plagas. Posteriormente, TensorFlow Lite permitirá que la Raspberry Pi procese las fotografías en tiempo real o de manera local, sin necesidad de depender completamente de servicios en la nube.

El sistema analiza esas imágenes y detecta automáticamente si hay señales de enfermedades, plagas o cualquier anomalía visible. En cuanto encuentra algo sospechoso, manda una notificación directamente al teléfono del usuario a través de un bot de Telegram, junto con la foto y la ubicación aproximada del área afectada. Así, el usuario puede reaccionar rápido antes de que se propague más la plaga o enfermedad a más zonas.

El proyecto integra tecnologías como dron de monitoreo aéreo, Raspberry Pi~4, inteligencia artificial, procesamiento digital de imágenes, comunicación inalámbrica (WiFi) y APIs de Telegram. La combinación de estas herramientas busca ofrecer una alternativa moderna y accesible para apoyar las tareas de supervisión agrícola, especialmente en cultivos donde el monitoreo constante resulta complicado.

La finalidad principal del sistema es contribuir a la detección temprana de enfermedades y plagas en cultivos de jitomate, reduciendo tiempos de supervisión y facilitando el monitoreo de grandes extensiones de terreno. Además, se busca disminuir pérdidas económicas, optimizar recursos y apoyar la modernización tecnológica dentro del sector agrícola.

\thispagestyle{empty}
